\documentclass[11pt,a4paper,fontset=fandol]{ctexart}

\usepackage[a4paper,top=2.25cm,bottom=2.45cm,left=2.25cm,right=2.25cm,headheight=18pt,headsep=0.55cm,footskip=24pt]{geometry}
\usepackage{amsmath,amssymb,amsthm}
\usepackage{fancyhdr}
\usepackage{lastpage}
\usepackage{enumitem}
\usepackage{titlesec}
\usepackage{xcolor}
\usepackage{hyperref}
\usepackage{bookmark}
\usepackage[most]{tcolorbox}
\usepackage{setspace}
\usepackage{array}

\hypersetup{
  colorlinks=true,
  linkcolor=blue!50!black,
  urlcolor=blue!50!black
}

\setstretch{1.13}
\setlength{\parindent}{2em}
\setlength{\parskip}{0.25em}
\allowdisplaybreaks
\setlist[enumerate]{itemsep=0.45em, topsep=0.35em, leftmargin=2.4em}
\setlist[itemize]{itemsep=0.25em, topsep=0.25em, leftmargin=1.9em}

\definecolor{mainblue}{RGB}{36,83,140}
\definecolor{lightblue}{RGB}{241,246,252}
\definecolor{lightgray}{RGB}{246,246,246}
\definecolor{rulegray}{RGB}{110,118,128}

\titleformat{\section}
  {\Large\bfseries\color{mainblue}}
  {\thesection.}
  {0.45em}
  {}
  [\vspace{0.2em}\color{rulegray}\titlerule]
\titlespacing*{\section}{0pt}{1.1em}{0.7em}

\pagestyle{fancy}
\fancyhf{}
\fancyhead[L]{\small 递推计数周测 B 卷}
\fancyhead[C]{\small 组合专题：状态设计与综合变式}
\fancyhead[R]{\small 刘文博}
\fancyfoot[C]{\small 第 \thepage 页 / 共 \pageref{LastPage} 页}
\renewcommand{\headrulewidth}{0.55pt}
\renewcommand{\footrulewidth}{0.35pt}

\newtcolorbox{infobox}[1]{
  breakable,
  colback=lightblue,
  colframe=mainblue,
  boxrule=0.7pt,
  arc=1mm,
  title=\textbf{#1},
  fonttitle=\bfseries
}

\newenvironment{solution}{\par\noindent\textbf{参考解答：}}{\hfill$\square$\par}
\newcommand{\answerspace}{\vspace{2.3cm}}
\newcommand{\biganswerspace}{\vspace{3.2cm}}

\begin{document}

\thispagestyle{empty}
\begin{center}
{\fontsize{22}{28}\selectfont\bfseries 递推计数周测 B 卷\par}
\vspace{0.4cm}
{\large 适合小学高年级至初中低年级数学竞赛训练\par}
\end{center}

\begin{infobox}{考试说明}
\begin{itemize}
  \item 测试时间：75 分钟
  \item 满分：100 分
  \item B 卷更强调状态设计与模型迁移，请特别注意：分类是否完整、是否重叠、初始条件是否写对。
\end{itemize}
\end{infobox}

\section{试题}

\begin{enumerate}
  \item （12 分）每次可以上 1 级或 2 级台阶，但不能连续两次都走 2 级。求上到第 8 级台阶的走法数。
  \answerspace

  \item （12 分）长度为 6 的 0,1 串中，1 的个数为偶数。这样的串有多少个？
  \answerspace

  \item （12 分）长度为 7 的 0,1 串中，不允许出现相邻两个 1，并且最后一位必须是 1。这样的串有多少个？
  \answerspace

  \item （12 分）用 $2\times 1$ 骨牌和 $2\times 2$ 方砖铺满 $2\times 6$ 长方形，有多少种铺法？
  \biganswerspace

  \item （12 分）证明：长度为 $n$ 的 0,1 串中，不允许出现相邻两个 1 的串数，与“每次上 1 级或 2 级台阶，上到第 $n+1$ 级”的走法数相同。
  \biganswerspace

  \item （12 分）长度为 9 的 A、B 串中，不允许出现连续三个 A。这样的串有多少个？
  \answerspace

  \item （14 分）已知错位排列满足
  \[
  D_n=(n-1)(D_{n-1}+D_{n-2}).
  \]
  用这个递推公式求 $D_7$，并解释它为什么也属于递推计数问题。
  \biganswerspace

  \item （14 分）用 $1\times 1$、$1\times 2$、$1\times 3$ 三种砖铺满 $1\times 7$ 长条，有多少种铺法？
  \biganswerspace
\end{enumerate}

\newpage
\section{详细解答}

\begin{enumerate}
  \item \begin{solution}
  设合法走法数为 $f_n$。

  若最后一步是 1，则前面有 $f_{n-1}$ 种；若最后一步是 2，由于不能连续两次走 2，末尾必须看成“12”结构，所以前面有 $f_{n-3}$ 种。

  因此
  \[
  f_n=f_{n-1}+f_{n-3} \qquad (n\ge 4).
  \]

  初始条件：
  \[
  f_1=1,\qquad f_2=2,\qquad f_3=3.
  \]

  于是
  \[
  f_4=4,\quad f_5=6,\quad f_6=9,\quad f_7=13,\quad f_8=19.
  \]

  所以答案为
  \[
  19.
  \]
  \end{solution}

  \item \begin{solution}
  设
  \[
  x_n=\text{长度为 $n$ 的 0,1 串中，1 的个数为偶数的个数},
  \]
  \[
  y_n=\text{长度为 $n$ 的 0,1 串中，1 的个数为奇数的个数}.
  \]

  在长度为 $n-1$ 的串后添 0，不改变奇偶性；添 1，会改变奇偶性，所以
  \[
  x_n=x_{n-1}+y_{n-1},
  \qquad
  y_n=x_{n-1}+y_{n-1}.
  \]

  这说明偶数个 1 和奇数个 1 的数量始终相同，而长度为 6 的 0,1 串总共有
  \[
  2^6=64
  \]
  个，所以偶数个 1 的有
  \[
  64\div 2=32
  \]
  个。
  \end{solution}

  \item \begin{solution}
  设长度为 $n$ 的不含相邻两个 1 的 0,1 串个数为 $a_n$，则
  \[
  a_n=a_{n-1}+a_{n-2},
  \qquad
  a_1=2,
  \quad
  a_2=3.
  \]

  所以
  \[
  a_3=5,\quad a_4=8,\quad a_5=13.
  \]

  现在要求长度为 7 的合法串且最后一位必须是 1，那么倒数第二位必定是 0，所以末尾固定为“01”，前面前 5 位只需任意合法。

  因此答案就是
  \[
  a_5=13.
  \]
  \end{solution}

  \item \begin{solution}
  设铺法数为 $s_n$。

  看最左端：
  \begin{itemize}
    \item 若先放一块竖着的 $2\times 1$ 骨牌，则剩下有 $s_{n-1}$ 种；
    \item 若先铺满一个 $2\times 2$ 小块，则可以用两块横骨牌，或用一块 $2\times 2$ 方砖，所以这一类共有 $2s_{n-2}$ 种。
  \end{itemize}

  因此
  \[
  s_n=s_{n-1}+2s_{n-2}.
  \]

  初始条件为
  \[
  s_1=1,\qquad s_2=3.
  \]

  所以
  \[
  s_3=5,\quad s_4=11,\quad s_5=21,\quad s_6=43.
  \]

  因此答案为
  \[
  43.
  \]
  \end{solution}

  \item \begin{solution}
  设
  \[
  a_n=\text{长度为 $n$ 的 0,1 串中，不允许出现相邻两个 1 的串数},
  \]
  \[
  f_n=\text{每次上 1 级或 2 级台阶，上到第 $n$ 级的走法数}.
  \]

  我们已经知道：
  \[
  a_n=a_{n-1}+a_{n-2},
  \qquad a_1=2,\ a_2=3;
  \]
  \[
  f_n=f_{n-1}+f_{n-2},
  \qquad f_1=1,\ f_2=2.
  \]

  观察初值：
  \[
  a_1=f_2,
  \qquad
  a_2=f_3.
  \]

  又因为它们满足完全相同的递推关系，所以对所有 $n$ 都有
  \[
  a_n=f_{n+1}.
  \]

  这就证明了两类问题的答案相同。
  \end{solution}

  \item \begin{solution}
  设合法串数为 $b_n$。

  结尾只能分成三类：B、AB、AAB，因此
  \[
  b_n=b_{n-1}+b_{n-2}+b_{n-3}.
  \]

  初始条件：
  \[
  b_1=2,\qquad b_2=4,\qquad b_3=7.
  \]

  所以
  \[
  b_4=13,\quad b_5=24,\quad b_6=44,
  \]
  \[
  b_7=81,\quad b_8=149,\quad b_9=274.
  \]

  因此答案为
  \[
  274.
  \]
  \end{solution}

  \item \begin{solution}
  先用递推求值。

  已知
  \[
  D_1=0,\qquad D_2=1.
  \]

  依次得到
  \[
  D_3=2(D_2+D_1)=2,
  \]
  \[
  D_4=3(D_3+D_2)=9,
  \]
  \[
  D_5=4(D_4+D_3)=44,
  \]
  \[
  D_6=5(D_5+D_4)=265,
  \]
  \[
  D_7=6(D_6+D_5)=6(265+44)=1854.
  \]

  所以
  \[
  D_7=1854.
  \]

  为什么这也属于递推计数？因为它本质上也是把“规模为 $n$ 的问题”拆成“规模为 $n-1$ 和 $n-2$ 的问题”。

  也就是说，我们并不是直接去数所有 $n$ 个对象的错位排列，而是通过分析对象 1 的去向，把问题转化成更小规模的两个子问题，再写成递推式。这正是递推计数的核心思想。
  \end{solution}

  \item \begin{solution}
  设铺法数为 $t_n$。

  看最后一块：
  \begin{itemize}
    \item 若最后一块是 $1\times 1$，前面有 $t_{n-1}$ 种；
    \item 若最后一块是 $1\times 2$，前面有 $t_{n-2}$ 种；
    \item 若最后一块是 $1\times 3$，前面有 $t_{n-3}$ 种。
  \end{itemize}

  因此
  \[
  t_n=t_{n-1}+t_{n-2}+t_{n-3}.
  \]

  初始条件：
  \[
  t_1=1,\qquad t_2=2,\qquad t_3=4.
  \]

  递推得到
  \[
  t_4=7,\quad t_5=13,\quad t_6=24,\quad t_7=44.
  \]

  所以答案为
  \[
  44.
  \]
  \end{solution}
\end{enumerate}

\end{document}
